\subsection{Simulation study}
\label{subsec:sim-study}

\begin{figure}[htbp]
    \centering
    \includegraphics[width=1\textwidth]{figures/scenarios.pdf}
    \caption{
        Simulation scenarios used for benchmarking. Each panel corresponds to a specific scenario and shows the relationship between target phylodynamic rates and their covariates.
        (a) \texttt{Epi-1}: constant reproduction number $R_t$ over time;
        (b) \texttt{Epi-2}: linearly decreasing reproduction number $R_t$ over time;
        (c) \texttt{Epi-3}: non-monotonic reproduction number $R_t$ over time;
        (d) \texttt{Epi-4}: migration rates $m_{ij}$ between five populations as a function of a continuous covariate $x_{ij}$ defined for each ordered pair of populations $(i,j)$;
        (e) \texttt{FBD-1}: constant speciation and extinction rates $\lambda$ and $\mu$ over time;
        (f) \texttt{FBD-2}: linearly varying speciation and extinction rates $\lambda$ and $\mu$ over time;
        (g) \texttt{FBD-3}: non-linear speciation rate $\lambda$ and a pulse-shaped extinction rate $\mu$ over time;
        (h) \texttt{FBD-4}: time-varying speciation and extinction rates under combinations of two binary traits $\mathcal{S} = \{00, 01, 10, 11\}$, where rates are denoted by $\lambda^{(s)}$ and $\mu^{(s)}$, for each state $s \in \mathcal{S}$. Speciation rates are linear in time, while extinction rates follow a pulse-shaped trajectory. Only the first trait influences diversification, determining whether $\lambda^{(s)}$ increases or decreases over time and controlling the magnitude of the pulse in $\mu^{(s)}$.
    }
    \label{fig:scenarios}
\end{figure}

We benchmarked our approach on eight simulation scenarios, grouped into two classes of multi-type birth--death models~\cite{stadler2013uncovering}: epidemiological processes (with removal after sampling) and fossilized birth--death (FBD) processes~\cite{Heath2014FBD} (with ancestor sampling and complete present-day sampling) (Fig.~\ref{fig:scenarios}).
Across these scenarios, we test different functional relationships between predictors and target phylodynamic rates, ranging from constant, linear, nonlinear, to interacting effects. 
For each scenario, in addition to the covariates that directly drive the phylodynamic rates of interest, we include additional random predictors that do not affect the evolutionary dynamics, in order to assess robustness to irrelevant covariates and evaluate potential overfitting.

All epidemiological scenarios share the same become-uninfectious rate $\delta = \mu + \psi = 0.07$, the same sampling probability upon removal $s = \frac{\psi}{\mu + \psi} = 0.15$, and the same simulation duration $T_{\text{EPI}} = 250$, where $\psi$ denotes the sampling rate through time and $\mu$ the death rate. All FBD scenarios assume a constant sampling rate $\psi = 0.2$ and a simulation duration $T_{\text{FBD}} = 35$.

\subsubsection{Scenarios \texttt{Epi-1} to \texttt{Epi-3}: time-varying reproduction number}

Scenarios \texttt{Epi-1}, \texttt{Epi-2}, and \texttt{Epi-3} consider a single homogeneous population, where the inference task is to estimate a time-varying reproduction number defined as
\[
R_t = \frac{\lambda}{\mu + \psi}.
\]

The function $R_t$ is represented as a piecewise constant vector $\mathbf{R} = (R_1, \dots, R_{10})$, defined over ten equal-length time intervals of duration $T_{\text{EPI}} / 10 = 25$ time units each.

In \texttt{Epi-1}, $R_t$ is constant over time, with $R_i = 1.2$ for all $i = 1, \dots, 10$. In \texttt{Epi-2}, $R_t$ decreases linearly across time bins, from $R_1 = 1.5$ to $R_{10} = 1.0$. In \texttt{Epi-3}, $R_t$ first increases linearly from $R_1 = 1.2$ to $R_5 = 1.5$, and then decreases linearly from $R_6 = 1.5$ to $R_{10} = 1.0$.

The predictor set consists of a time vector $\mathbf{t} = (t_1, \dots, t_{10})$ with $t_i = i$, and a stochastic covariate vector $\mathbf{x} = (x_1, \dots, x_{10})$ generated via a standard normal random walk over the same time interval.
In this setting, BELLA takes as input pairs $(t_i, x_i)$ and outputs the corresponding reproduction number $R_i$.

\subsubsection{Scenario \texttt{Epi-4}: pairwise migration rates in a structured population}

Scenario \texttt{Epi-4} consists of five populations with reproduction numbers $(0.8, 1.0, 1.2, 1.4, 1.6)$, respectively. The inference task is to estimate the $5 \times 4 = 20$ pairwise migration rates $m_{ij}$ between populations.

The predictor set includes a covariate $x_{ij}$ defined for each ordered population pair $(i,j)$ with $i \neq j$. For each pair, $x_{ij}$ is sampled independently from a uniform distribution,
\[
x_{ij} \sim \mathcal{U}(-1,1).
\]

These predictors are mapped to the target migration rates using a logistic transformation,
\[
m_{ij} = \frac{0.04}{1 + \exp(- 8 x_{ij})}.
\]

In addition, a random covariate matrix $\mathbf{z} = (z_{ij})$ is generated by independently sampling each entry as
\[
z_{ij} \sim \mathcal{U}(-1,1),
\]
for all population pairs $(i,j)$.

In this setting, BELLA takes as input feature pairs $(x_{ij}, z_{ij})$ and outputs the corresponding migration rate $m_{ij}$ for each population pair.

\subsubsection{Scenarios \texttt{FBD-1} to \texttt{FBD-3}: time-varying speciation and extinction rates}

Scenarios \texttt{FBD-1}, \texttt{FBD-2}, and \texttt{FBD-3} consider a single homogeneous population, where the inference task is to estimate time-varying speciation and extinction rates, denoted by $\boldsymbol{\lambda} = (\lambda_1, \dots, \lambda_{10})$ and $\boldsymbol{\mu} = (\mu_1, \dots, \mu_{10})$.

Both rates are represented as piecewise constant vectors over ten equal-length time intervals of duration $T_{\text{FBD}} / 10 = 3.5$ time units each.

In \texttt{FBD-1}, speciation and extinction rates are constant over time, with $\lambda_i = 0.2$ and $\mu_i = 0.1$ for all $i = 1, \dots, 10$. In \texttt{FBD-2}, the speciation rate decreases linearly from $0.4$ to $0.1$, while the extinction rate increases linearly from $0.1$ to $0.2$ across time bins. In \texttt{FBD-3}, the speciation rate follows a non-linear decreasing pattern given by
\[
\lambda_i = 
\begin{cases}
0.4, & \text{if } 1 \le i \le 5, \\
0.1 & \text{if } 6 \le i \le 8, \\ 
0.01 & \text{if } 9 \le i \le 10,
\end{cases}
\]
with higher values in the first five time intervals, followed by a lower plateau and a final drop.

The extinction rate is mostly low but exhibits a sharp spike in the second half of the simulation, given by
\[
\mu_i = 
\begin{cases}
0.05, & \text{if } 1 \le i \le 7, \\
0.3 & \text{if } i = 8, \\ 
0.01 & \text{if } 9 \le i \le 10,
\end{cases}
\]
where a single elevated value occurs at the eighth time interval, followed by a return to low rates.

The predictor set consists of a time vector $\mathbf{t} = (t_1, \dots, t_{10})$ with $t_i = i$, and a stochastic covariate vector $\mathbf{x} = (x_1, \dots, x_{10})$ generated via a standard normal random walk over the same time interval.
In this setting, BELLA takes as input pairs $(t_i, x_i)$ and outputs the corresponding rates $(\lambda_i, \mu_i)$.

\subsubsection{Scenario \texttt{FBD-4}: time-varying and trait-dependent speciation and extinction rates}

Scenario \texttt{FBD-4} considers a structured population with four discrete states corresponding to combinations of two binary traits, $\mathcal{S} = \{00, 01, 10, 11\}$. We allow migration between states corresponding to lineage transitions in which a single trait changes at a time. Specifically, migration is restricted to adjacent states in the Hamming space, with all permitted transitions occurring at rate $m = 0.1$. The inference task is to estimate time-varying speciation and extinction rates associated with each state.

Both rates are represented as state-dependent, piecewise constant vectors over ten equal-length time intervals of duration $T_{\text{FBD}} / 10 = 3.5$ time units each: $\boldsymbol{\lambda}^{(s)} = (\lambda^{(s)}_1, \dots, \lambda^{(s)}_{10})$ and $\boldsymbol{\mu}^{(s)} = (\mu^{(s)}_1, \dots, \mu^{(s)}_{10})$ for each state $s \in \mathcal{S}$.

Only the first trait affects diversification. For states with the first trait unset ($s \in \{00,01\}$), the speciation rate decreases linearly,
\[
\boldsymbol{\lambda}^{(00)} = \boldsymbol{\lambda}^{(01)} = (0.6, \dots, 0.1),
\]
whereas for states with the first trait set ($s \in \{10,11\}$), the speciation rate increases linearly,
\[
\boldsymbol{\lambda}^{(10)} = \boldsymbol{\lambda}^{(11)} = (0.2, \dots, 0.5).
\]

The extinction rate exhibits a pulse-shaped pattern. For states with the first trait unset,
\[
\mu^{(00)}_i = \mu^{(01)}_i
\begin{cases}
1.0, & \text{if } 5 \le i \le 8, \\
0.1, & \text{otherwise},
\end{cases}
\]
while for states with the first trait set,
\[
\mu^{(10)}_i = \mu^{(11)}_i
\begin{cases}
0.5, & \text{if } 5 \le i \le 8, \\
0.1, & \text{otherwise},
\end{cases}
\]
Thus, the first trait determines both the direction of the speciation trend and the magnitude of the extinction pulse, while the second trait has no effect on diversification.

The predictor set consists of a time covariate $\mathbf{t} = (t_1, \dots, t_{10})$ with $t_i = i$, a binary covariate $\mathbf{d} = (d_1, \dots, d_{10})$ encoding the presence of an extinction spike, defined as
\[
d_i =
\begin{cases}
1, & \text{if } 5 \le i \le 8, \\
0, & \text{otherwise},
\end{cases}
\]
and binary indicators encoding the trait values for each state.

In this setting, BELLA takes as input tuples $(t_i, d_i, \tau^{(1)}_s, \tau^{(2)}_s)$, where $\tau^{(1)}_s$ and $\tau^{(2)}_s$ denote the values of the two traits for state $s$, and outputs the corresponding rates $(\lambda^{(s)}_i, \mu^{(s)}_i)$.

\subsubsection{Inference settings}

For each scenario, we simulated 100 phylogenetic trees with between 200 and 500 tips using the Phylogenie Python package (v3.11). We performed MCMC inference on each tree using the BDMM-Prime package~\cite{Vaughan2025BDMM} implemented in BEAST 2~\cite{bouckaert2019beast} coupled with different inference approaches: (i) a predictor-agnostic (PA) model that estimates each parameter independently, without conditioning on covariates; (ii) a GLM implementation~\cite{cecivale_glmprior_2025} that estimates an intercept and growth factors defining an exponential mapping between predictors and phylodynamic parameters, and (iii) nine different BELLA configurations (Tab.~\ref{tab:models}).
For analyses requiring multiple parameters (e.g., speciation and extinction rates in FBD scenarios), a separate network was used for each target parameter, with the output layer consisting of a single neuron.

For the predictor-agnostic model and GLMs, we employed uniform priors for the phylodynamic parameters: \( \mathcal{U}(1, 5) \) for reproduction numbers,  \( \mathcal{U}(0, 0.05) \) for migration rates, and \( \mathcal{U}(0, 2) \) for speciation and extinction rates. 
We used standard normal priors \(\mathcal{N}(0, 1)\) for the intercept and growth factors in the GLMs. 
For BELLA, we constrained the output phylodynamic parameters through a logistic output activation to match the bounds of the corresponding uniform priors.
All MCMC runs employed chains of 10,000,000 steps. 
For all analyses, the first 10\% of samples were discarded as burn-in, and convergence was assessed by verifying that the posterior effective sample size exceeded 200.
